\documentclass[11pt]{article}
\usepackage[T1]{fontenc}
\usepackage{lmodern,amsmath,amssymb,amsthm,booktabs,array,longtable}
\usepackage[margin=1in]{geometry}
\usepackage[colorlinks=true,urlcolor=blue,citecolor=blue,linkcolor=blue]{hyperref}
\newtheorem{theorem}{Theorem}
\newtheorem{corollary}[theorem]{Corollary}
\newcommand{\tr}{\operatorname{tr}}
\newcommand{\eps}{\varepsilon}
\title{Spectral extensions of Vinay (2019) and Wicklin (2021):\\
resolving Spezia's conjectures in OEIS A322844, COMMENTS,\\
28 December 2018 and 6 January 2019}
\author{Haobo Ma\\\small ChronoAI Pte Ltd, Singapore
\and Wenlin Zhang\\\small National University of Singapore\\
\small ChronoAI Pte Ltd, Singapore}
\date{8 November 2026}
\begin{document}
\maketitle
\begin{abstract}
For the conjectures in OEIS A322844, COMMENTS, first posed on
28 December 2018 and 6 January 2019 \cite{posing}, we extend the arithmetic-row rank
argument of Vinay (2019) and Wicklin (2021) to an explicit spectral formula.
For the $n\times n$
matrix containing $1,\ldots,n^2$ in alternating row directions, a rank-two
factorization gives the signed second exterior trace and the vanishing of
every higher exterior trace. The same calculation gives the two nonzero
eigenvalues and their power-sum recurrence, including the even-dimensional
eigenvalue interpretation of A317614 and the conjecture in A317297. We derive the linked sum
and permanent formulas, the generating functions and reciprocal sums,
and distinguish three displayed exponential functions from their usual
exponential-generating-function interpretations.
\end{abstract}
\noindent\textbf{License.} Creative Commons Attribution 4.0 International (CC BY 4.0).

\section{The two conjectures and the characteristic polynomial}
Under the heading ``Conjectures'', OEIS A322844, COMMENTS \cite{posing},
asserts that for $n>1$ its sequence is the ``absolute value of the trace of
the 2nd exterior power'' of a serpentine matrix, and that ``For $k>2$''
the higher exterior trace is ``equal to zero''.
The first formula appears in Spezia's original submission, revision 2
(28 December 2018); the higher-exterior assertion was added in revision 5
(6 January 2019). Revision 47 approved the entry on 9 February 2019.
We prove both by writing the matrix as a product of an
$n\times2$ matrix and a $2\times n$ matrix. The result also determines its
characteristic polynomial, extending the determinant formula recorded in
OEIS A317614 \cite{matrix}. Vinay's column-span argument (2019), followed
by Wicklin's exposition (2021), already proves that an arithmetic-row
matrix has rank at most two \cite{vinay,wicklin}. It directly implies
higher-exterior vanishing for the present matrix and for arbitrary row
directions; the explicit quadratic and signed second coefficient below
are the additional calculation. The even-order construction and its sum
identities occur earlier in De Los Reyes et al. (2007) \cite{rectangles}.
The 2012 Miranda construction announcements \cite{miranda2012} provide
further antecedents; the relation of their unavailable mathematical texts
to this spectral coefficient remains unresolved. None of these comparisons
is a priority claim. Throughout, rows and columns are numbered from $1$.
Define
\begin{equation}\label{matrix}
 M(n)_{ij}=\begin{cases}
 j+n(i-1),&i\text{ odd},\\
 ni-j+1,&i\text{ even},
 \end{cases}\qquad 1\leq i,j\leq n.
\end{equation}
For a square matrix $A$, put $e_k(A)=\tr(\bigwedge^k A)$, with
$e_0(A)=1$ and $e_k(A)=0$ for $k>\dim A$. In the standard exterior basis,
the diagonal entries of $\bigwedge^k A$ are the principal $k\times k$
minors. Expanding the determinant by choosing either $x$ or an entry of
$-A$ in each row therefore gives
\begin{equation}\label{exterior}
 \det(xI-A)=\sum_{k=0}^{n}(-1)^k e_k(A)x^{n-k}.
\end{equation}
This is the characteristic-coefficient convention in the posing entry.

\begin{theorem}\label{main}
Write $\eps=n\bmod2$ and
\begin{equation}\label{constants}
 T_n=\frac{n(n^2+\eps)}2,\qquad
 A_n=\frac{n^2}{12}\{3(1+n^2)-2(2+n^2)\eps\}
 =\begin{cases}
 n^2(n^2+1)/4,&n\text{ even},\\
 n^2(n^2-1)/12,&n\text{ odd}.
 \end{cases}
\end{equation}
For every integer $n\geq2$,
\begin{equation}\label{charpoly}
 \det(xI-M(n))=x^{n-2}(x^2-T_nx-A_n).
\end{equation}
Consequently $e_1(M(n))=T_n$, $e_2(M(n))=-A_n$, and
$|e_2(M(n))|=A_n=a(n)$ in A322844. For every $n\geq0$ and $k>2$,
$e_k(M(n))=0$.
\end{theorem}

\section{A two-dimensional calculation}
For $1\leq i,j\leq n$, set
\[
 s_i=(-1)^{i-1},\qquad
 w_j=j-\frac{n+1}{2},\qquad
 u_i=ni-\frac{n-1}{2}.
\]
Direct substitution in the two cases of \eqref{matrix} gives
\begin{equation}\label{factorization}
 M(n)_{ij}=u_i+s_iw_j,\qquad M(n)=UV,\qquad
 U=\begin{pmatrix}u_1&s_1\\\vdots&\vdots\\u_n&s_n\end{pmatrix},\quad
 V=\begin{pmatrix}1&\cdots&1\\w_1&\cdots&w_n\end{pmatrix}.
\end{equation}
The centered sums are
\begin{equation}\label{moments}
 \sum_i w_i=0,\quad
 \sum_i w_i^2=\frac{n(n^2-1)}{12},\quad
 \sum_i s_i=\eps,\quad
 \sum_i i s_i=\frac{-n+(2n+1)\eps}{2}.
\end{equation}
The first two follow from $\sum i=n(n+1)/2$ and
$\sum i^2=n(n+1)(2n+1)/6$, by expansion. Pairing consecutive terms
gives the last two: each pair contributes $0$ and $-1$, respectively;
when $n$ is odd, the remaining term contributes $1$ and $n$.
It follows that $\sum_i s_iw_i=-n(1-\eps)/2$.
Also $u_i=nw_i+(n^2+1)/2$, so
\begin{equation}\label{small}
 B:=VU=\begin{pmatrix}
 n(n^2+1)/2&\eps\\
 n^2(n^2-1)/12&-n(1-\eps)/2
 \end{pmatrix}.
\end{equation}
Its trace is $T_n$, and its determinant is
\[
 -\frac{n^2(n^2+1)(1-\eps)}4
 -\frac{\eps n^2(n^2-1)}{12}=-A_n.
\]

For $n\geq2$ the columns of $U$ are linearly independent: their first
two coordinates form a $2\times2$ matrix of determinant
$-u_1-u_2=-(2n+1)\ne0$. The image of $U$ is invariant under $M$, since
$MU=UB$, and the induced map on the quotient by this image is zero,
since $M=UV$. Extend the two columns of $U$ to a basis. In that basis
$M$ has block form $\left(\begin{smallmatrix}B&C\\0&0\end{smallmatrix}\right)$,
which proves \eqref{charpoly} by taking the determinant of $xI-M$.
Comparison with \eqref{exterior} proves the asserted exterior traces.
Both parity formulas for $A_n$ are positive when $n\geq2$, so its
absolute-value assertion follows as well. For $n=0$ the matrix is empty;
for $n=1$ it is $(1)$. Their higher exterior spaces are zero, which
completes the boundary cases.\hfill$\square$

\begin{corollary}[Arbitrary row directions; not posed in the source]\label{directions}
Choose any signs $s_i\in\{-1,1\}$ to select forward or backward rows,
and retain $u_i,w_j$ from \eqref{factorization}. Put
$S=\sum_i s_i$, $W=\sum_i s_iw_i$,
$C=n(n^2+1)/2$, and $D=n^2(n^2-1)/12$.
For $n\geq2$, the resulting matrix has characteristic polynomial
\[
 x^{n-2}\{x^2-(C+W)x+CW-SD\}.
\]
In particular every exterior trace of order above two vanishes for any
row-direction pattern.
\end{corollary}
\begin{proof}
The same factorization gives $VU=\left(\begin{smallmatrix}C&S\\D&W\end{smallmatrix}\right)$.
The two columns of $U$ are independent because $u_1,u_2$ are positive and
unequal and $s_1,s_2$ are signs. Apply the same invariant-subspace argument.
\end{proof}

\section{Spectral consequences and the linked matrix sequences}
\begin{corollary}\label{spectrum}
For $n\geq2$, $M(n)$ has rank two and is diagonalizable over $\mathbb R$.
Its two nonzero eigenvalues are
\[
 \lambda_{\pm}=\frac{T_n\pm\sqrt{T_n^2+4A_n}}2,
 \qquad \lambda_-<0<\lambda_+,
\]
and its other $n-2$ eigenvalues are zero. If $p_r=\tr(M(n)^r)$ for
$r\geq1$, and $p_0=2$, then
\[
 p_1=T_n,\qquad p_2=T_n^2+2A_n,\qquad
 p_r=T_np_{r-1}+A_np_{r-2}\quad(r\geq2).
\]
Here $p_0$ counts the two nonzero eigenvalues, not the trace of $I_n$.
\end{corollary}
\begin{proof}
Since $\det B=-A_n\ne0$, the rank of $UV$ is two. The roots of its
quadratic factor are distinct and nonzero, with negative product.
The kernel of $M$ has dimension $n-2$, equal to the multiplicity of zero
in \eqref{charpoly}; together with the other two eigenspaces this gives
a basis of eigenvectors. Finally $p_r=\lambda_+^r+\lambda_-^r$, and each
root satisfies $\lambda^2=T_n\lambda+A_n$. Summing the resulting
recurrence proves the formulas, including the stated convention at $r=0$.
\end{proof}
At $n=2$, the matrix is $\left(\begin{smallmatrix}1&2\\4&3\end{smallmatrix}\right)$
and its polynomial is $x^2-4x-5$. At $n=3$ it is $x(x^2-15x-6)$.
For $n\geq3$, \eqref{charpoly} gives $\det M(n)=0$, recovering the
determinant identity in A317614; its complete initial values are $1,1,-5$
at $n=0,1,2$. The determinant is already labelled proved in revision 89
of A317614 (24 August 2018), without an accompanying argument.
The rank pattern $0,1,2,2,\ldots$ was added there as the A130130
cross-reference on 9 April 2023, revision 151 \cite{matrix}.

\begin{corollary}[Even-dimensional eigenvalues]\label{even}
For $n=2m\geq2$, the nonzero eigenvalues are
\[
 \lambda_+=m(4m^2+1),\qquad \lambda_-=-m.
\]
In particular $\rho(M(2m))=m(4m^2+1)$, and the minimum eigenvalue
is $-m$. This gives a spectral interpretation of A317614, FORMULA,
28 December 2018, and the size-indexed version of Spezia's conjecture
in A317297, COMMENTS, 27 December 2018 \cite{eigen}.
\end{corollary}
\begin{proof}
When $n$ is even, $B$ in \eqref{small} is lower triangular. Its diagonal
entries are $n(n^2+1)/2$ and $-n/2$, which proves the assertion.
\end{proof}
The two entries use a separate parameter for the dimension.
In terms of the actual dimension, the A317297 sequence is
$b(r)=(r-1)(4(r-1)^2+1)=\rho(M(2r-2))$ for $r\geq2$;
$b(1)=0$ agrees with the zero-dimensional linear map.
There is a literal indexing discrepancy in its displayed matrix definition:
the stated dimension is $2(r-1)$, but the row-step coefficient is printed
as $r$. At $r=3$ that displayed $4\times4$ matrix is
\[
 \begin{pmatrix}1&2&3&4\\6&5&4&3\\7&8&9&10\\12&11&10&9\end{pmatrix},
 \qquad \det(xI-N)=x^2(x-26)(x+2),
\]
whose maximum is $26$, not $b(3)=34$. Replacing the row-step coefficient
by the dimension gives $M(4)$ and restores the claimed sequence.
The minimum eigenvalue formula is a further corollary, not a conjecture
asserted in A317614. More generally, for any positive row-step $h$ and even
dimension $d$, the corresponding two eigenvalues are
$d(hd+1)/2$ and $-d/2$. Indeed use
$u_i=hi-(h-1)/2$, $w_j=j-(h+1)/2$ in \eqref{factorization};
then $VU$ is triangular because $\sum s_i=0$.

\begin{corollary}[The polynomial and triangle links in A317297]\label{triangle}
Put $b(r)=(r-1)(4(r-1)^2+1)$ for $r\geq1$. Its ordinary generating
function, with the source index, and its offset-normalized exponential
generating function are
\[
 \sum_{r\geq1}b(r)x^r=\frac{x^2(5+14x+5x^2)}{(1-x)^4},\qquad
 \sum_{m\geq0}b(m+1)\frac{x^m}{m!}=x(5+12x+4x^2)e^x.
\]
The recurrence is $b(r)=4b(r-1)-6b(r-2)+4b(r-3)-b(r-4)$ for $r\geq5$;
it also holds at $r=4$ with the polynomial extension $b(0)=-5$.
The further descriptions in A317297, COMMENTS and FORMULA (Oktar, 2024;
Navarrete, 2025), follow from
\[
 b(r)=\frac{(2r-1)^4-(2r-3)^4}{16}
     =3\binom{2r-1}{3}+2(r-1).
\]
For $r=1$ the binomial term is zero. The second identity also counts
$b(r)+1$ ternary words of length $L=2r-1$ with exactly one $1$ and
either zero or two $0$'s: choose the $1$ in $L$ ways and the two $0$'s,
if present, in $\binom{L-1}{2}$ ways. Thus the count is
$L+L\binom{L-1}{2}=b(r)+1$. Both polynomial identities follow by expansion.
If $c(j)$ denotes row $j$ of A246696 summed across the row (A246697),
numbered from $0$, then
\[
 c(2m)=4m^3+6m^2+5m+1,\qquad c(2m-1)=4m^3+m\ (m\geq1).
\]
Consequently $b(r)=c(2r-3)$ for $r\geq2$, proving the bisection
conjecture in A317297, CROSSREFS (entry by Pol, 1 September 2018) \cite{trianglelinks}.
These row-sum formulas also prove the equivalent closed-form conjectures
in A246697, FORMULA (Kimberling, 2014; Etienne, 2016; Spezia, 2021).
In full-index form, $c(j)=(2j^3+6j^2+9j+4+j(-1)^j)/4$; hence
\[
 \sum_{j\geq0}c(j)x^j=\frac{1+3x+5x^2+x^3+2x^4}{(1-x)^4(1+x)^2},\quad
 \sum_{j\geq0}c(j)\frac{x^j}{j!}
 =\frac{e^x(2x^3+12x^2+17x+4)-xe^{-x}}4.
\]
This proves its conjectured ordinary and exponential generating functions
(the latter equals the source's hyperbolic expression), and the recurrence
with signature $(2,1,-4,1,2,-1)$ for $j\geq6$.
\end{corollary}
\begin{proof}
The formulas for $b$ follow by applying
$4\Theta^3+\Theta$ to $(1-x)^{-1}$ or $e^x$, and taking fourth differences.
For the triangle, its defining recurrence makes successive entries of
each parity increase by $2$. Direct simultaneous induction on the row
gives first two entries $(2m^2+2m+1,2m^2+2)$ in even row $2m$ ($m\geq1$),
and $(2m^2-2m+2,2m^2+1)$ in odd row $2m-1$.
The initial rows are $(1)$ and $(2,3)$.
For $m\geq2$, from even row $2m-2$ to odd row $2m-1$, the last even-row
entry plus $2$ is the new second entry, and the penultimate entry plus $2$
is the new first entry. From odd row $2m-1$ to even row $2m$, the last
entry plus $2$ is the new first entry and the preceding entry plus $2$
is the new second entry. These are exactly the displayed polynomials.
Summing the two arithmetic progressions gives the formulas for $c$.
At row $0$, the even formula also gives $1$.
Applying its full-index polynomial in $\Theta$ to $(1-x)^{-1}$ and $e^x$
gives the displayed functions; multiplying by $(1-x)^4(1+x)^2$
gives the recurrence and its initial conditions.
\end{proof}

\begin{corollary}[The permutation and parity counts in A246696]\label{permutation}
For $m\geq1$, row $2m$ of A246696 contains $m$ even and $m+1$ odd entries, while
row $2m-1$ contains $m$ entries of each parity. Every positive integer
occurs exactly once. Thus the stated permutation property is true, but
the printed count $[(m+1)/2]$ of odd entries in row $2m$ is replaced
by $m+1$; already row $2$ is $(5,4,7)$.
\end{corollary}
\begin{proof}
Use the first two entries and the step-two progressions in the proof of
Corollary~\ref{triangle}. Row $2m$ starts with an odd entry and has length
$2m+1$; row $2m-1$ starts with an even entry and has length $2m$.
The two rows together contain exactly the even numbers from
$2m^2-2m+2$ to $2m^2+2m$, and exactly the odd numbers from $2m^2+1$
to $2m^2+4m+1$. These intervals are disjoint and successive paired
intervals abut within each parity: their next lower endpoints are the
current upper endpoints plus two. With the initial row $(1)$ this proves
the permutation assertion, including existence and uniqueness.
\end{proof}

The other matrix sequences cross-referenced by A317614 have short formulas
from the same decomposition. Put $\eps_m=m\bmod2$. For $n\geq1$,
\begin{align}
 R_i&:=\sum_j M(n)_{ij}=n^2i-\frac{n(n-1)}2,
 &C_j&:=\sum_i M(n)_{ij}=\frac{n(n^2+1)}2+
 \eps_n\left(j-\frac{n+1}2\right).\label{rowcol}
\end{align}
The even-order array already appears, transposed, in De Los Reyes,
Das, Midha and Vellaisamy (2007), Step A1; their Appendix (i)--(ii)
proves the corresponding sum identities \cite{rectangles}.
The formulas in \eqref{rowcol} give a uniform parity-dependent calculation.
These are respectively the row-sum triangle A241016 and the column-sum
triangle A317617 \cite{sums}: A241016 lists $R_1,\ldots,R_n$, whereas A317617 lists
$C_1,\ldots,C_n$.
The other matrix in A317614, COMMENTS (Barari, 31 October 2021;
Spezia, 5 November 2021), has entries
$H_{ij}=n-|i-(n+1)/2|-|j-(n+1)/2|$ and total sum $T_n$.
Indeed $\sum_i|i-(n+1)/2|$ is $n^2/4$ for even $n$ and $(n^2-1)/4$
for odd $n$, so $\sum_{ij}H_{ij}=n^3-2n\sum_i|i-(n+1)/2|=T_n$.
For $0\leq k<n$, the superdiagonal and subdiagonal sums are
\begin{align}
 D_k^+&:=\sum_{i=1}^{n-k}M(n)_{i,i+k}
       =\frac n2\big((n-k)^2+\eps_{n-k}\big),\label{upper}\\
 D_k^-&:=\sum_{i=k+1}^{n}M(n)_{i,i-k}
       =\frac{n(n^2-k^2+\eps_n+\eps_k)-2k\eps_n}{2}.\label{lower}
\end{align}
They give A325516 and A325655, respectively; for $n\geq2$, $D_1^+$ gives A323724 and
$D_1^-$ gives A323723 \cite{diagonals}. The $k=0$ instances both
give $T_n$. Formula \eqref{rowcol} follows by summing
\eqref{factorization}. To check \eqref{upper}--\eqref{lower}, define
\[
 H(m)=\sum_{i=1}^m i s_i=\frac{-m+(2m+1)\eps_m}{2}.
\]
For a diagonal with column $i+d$ and rows $a\leq i\leq b$, summation
of \eqref{factorization} gives
\begin{equation}\label{diagonalproof}
 n\frac{(a+b)(b-a+1)}2-\frac{(n-1)(b-a+1)}2
 +H(b)-H(a-1)+\left(d-\frac{n+1}2\right)(\eps_b-\eps_{a-1}).
\end{equation}
Use $(a,b,d)=(1,n-k,k)$ for \eqref{upper}, and
$(k+1,n,-k)$ for \eqref{lower}; expanding $H$ gives the stated formulas.

\begin{corollary}[First off-diagonal parity and terminal digits]\label{lastdigits}
For $n\geq2$, the first lower and upper diagonal sums in A323723 and A323724 are even.
The upper sum never has terminal digit $4$, as asserted in A323724,
COMMENTS \cite{diagonals}.
\end{corollary}
\begin{proof}
At $n=2m$, the lower and upper sums are respectively
$4m^3$ and $2m(2m^2-2m+1)$; at $n=2m+1$ they are
$2m(2m^2+3m+2)$ and $2m^2(2m+1)$, by \eqref{upper}--\eqref{lower}
with $k=1$. Each is twice an integer. For $m=0,1,2,3,4$ modulo $5$,
the upper sums have residues $(0,2,0,8,0)$ and $(0,6,0,6,8)$ modulo $10$,
respectively. Because each expression is twice an integer polynomial,
these five cases exhaust its terminal digits, none of which is $4$.
\end{proof}

The fixed-size antidiagonal is determined entrywise by
\begin{equation}\label{anti}
 M(n)_{i,n+1-i}=\begin{cases}(n-1)i+1,&i\text{ odd},\\
 (n+1)i-n,&i\text{ even}.
 \end{cases}
\end{equation}
This follows by substitution in \eqref{matrix}. Thus at $n=4$ it is
$(4,6,10,16)$. Although A317614 links A074147 under antidiagonals,
the fourth row of that triangle is $(6,8,10,12)$ \cite{antidiagonal};
that row is not the fixed-size antidiagonal of $M(4)$.
The cross-reference alone does not assert this fixed-size interpretation.

\begin{corollary}[Antidiagonal products]\label{antiproduct}
The sequence A382597 \cite{antiproductref} is
\[
 Q_n=\prod_{\substack{1\leq i\leq n\\ i\text{ odd}}}((n-1)i+1)
     \prod_{\substack{1\leq i\leq n\\ i\text{ even}}}((n+1)i-n).
\]
As recorded there, $Q_n$ has a trailing zero exactly when $n$ is even
and $n\notin\{0,2,6\}$.
\end{corollary}
\begin{proof}
The product follows from \eqref{anti}. For odd $n$ every factor is odd.
For even $n$ every factor is even. If $n\geq10$ and $n\not\equiv1\pmod5$,
one of the odd indices $1,3,5,7,9$ makes the first factor zero modulo $5$.
If $n\equiv1\pmod5$, one of $2,4,6,8,10$ makes the second factor zero
modulo $5$. The case $n=8$ has antidiagonal entry $10$;
$Q_0=1$, $Q_2=8$, $Q_4=3840$, and $Q_6=15814656$ finish the small cases.
\end{proof}

\begin{corollary}[The linked permanent]\label{permanent}
Let $\alpha_k=[z^k]\prod_{i=1}^n(u_i+s_i z)$ and
$\beta_k=[z^k]\prod_{j=1}^n(1+w_jz)$. Then the permanent in A322277
\cite{permanentref} is
\begin{equation}\label{permanentformula}
 \operatorname{per}M(n)=\sum_{k=0}^n k!(n-k)!\,\alpha_k\beta_k.
\end{equation}
\end{corollary}
\begin{proof}
Expand $\sum_\sigma\prod_i(u_i+s_iw_{\sigma(i)})$ by the set $S$ of rows
choosing the second summand. For $|S|=k$ and a fixed $k$-element image
$J=\sigma(S)$, exactly $k!(n-k)!$ permutations have that image.
The row products sum to $\alpha_k$ and the column products to $\beta_k$.
This also proves the empty-matrix value $1$.
\end{proof}

\section{Generating functions and reciprocal sums}\label{generating}
Write $a_n=A_n$, including $a_0=a_1=0$, and let
$\Theta=x\,d/dx$. The following gives coefficient-level meanings to all
the exponential expressions in A322844, FORMULA \cite{posing}.

\begin{corollary}\label{gf}
The ordinary generating function is
\[
 G(x)=\sum_{n\geq0}a_nx^n=
 \frac{x^2(5+6x+43x^2+20x^3+43x^4+6x^5+5x^6)}{(1-x^2)^5}.
\]
Its even and odd subsequence generating functions are respectively
\[
 \sum_{m\geq0}a_{2m}x^m=\frac{x(5+43x+43x^2+5x^3)}{(1-x)^5},\quad
 \sum_{m\geq0}a_{2m+1}x^m=\frac{2x(3+10x+3x^2)}{(1-x)^5}.
\]
Both subsequences satisfy the order-five recurrence with coefficients
$5,-10,10,-5,1$, starting at subsequence index $m=5$.
The exponential generating function in the usual convention is
\begin{equation}\label{trueegf}
 E(x):=\sum_{n\geq0}a_n\frac{x^n}{n!}
 =\frac{e^x(2x^4+12x^3+15x^2+3x)
       +e^{-x}(x^4-6x^3+9x^2-3x)}{12}.
\end{equation}
The usual exponential generating functions of $a_{2m}$ and $a_{2m+1}$
are $(4\Theta^4+\Theta^2)e^x$ and
$(4\Theta^4+8\Theta^3+5\Theta^2+\Theta)e^x/3$.
\end{corollary}
\begin{proof}
The parity formulas give $a_{2m}=4m^4+m^2$ and
$a_{2m+1}=(4m^4+8m^3+5m^2+m)/3$.
Apply these polynomials in $\Theta$ to $(1-x)^{-1}$ for the ordinary
functions, and to $e^x$ for their usual exponential counterparts.
The fifth backward difference of each degree-four polynomial vanishes,
giving the recurrence. For the full exponential function use
\[
 a_n=\frac{n^4}{6}+\frac{n^2}{12}
       +(-1)^n\left(\frac{n^4}{12}+\frac{n^2}{6}\right)
\]
and $\Theta^2e^x=(x^2+x)e^x$,
$\Theta^4e^x=(x^4+6x^3+7x^2+x)e^x$; replace $x$ by $-x$
in the alternating part. These calculations prove all identities.
\end{proof}
The overall expression labelled ``E.g.f.'' in A322844 is instead
\begin{equation}\label{shifted}
 F(x)=\sum_{n\geq0}\frac{a_nx^n}{(n+2)!}
 =\frac{e^{-x}(24+21x+9x^2+2x^3+x^4)-60
       +e^x(36-33x+15x^2-4x^3+2x^4)}{12x^2}.
\end{equation}
In particular $[x^2]F=5/24$, whereas $[x^2]E=5/2$.
The corrected usual EGF is \eqref{trueegf}. To verify the exact
interpretation \eqref{shifted}, differentiate $x^2F(x)$ twice:
the result is $E(x)$, with both integration constants zero at $0$.
The two displayed subsequence exponential expressions have the exact meanings
\begin{align*}
 H_o(x)&=\sum_{m\geq0}\frac{a_{2m+1}x^m}{(2m+1)!}
 =\frac{6x\cosh\sqrt{x}+\sqrt{x}(6+x)\sinh\sqrt{x}}{12},\\
 H_e(x)&=\sum_{m\geq0}\frac{a_{2m}x^m}{(2m)!}
 =\frac{x(8+x)\cosh\sqrt{x}+2\sqrt{x}(1+3x)\sinh\sqrt{x}}{4}.
\end{align*}
Indeed $H_e(t^2)=(E(t)+E(-t))/2$ and
$tH_o(t^2)=(E(t)-E(-t))/2$. These identities are formal power series,
so no square-root branch is involved. They distinguish the displayed
functions from subsequence EGFs with denominator $m!$.

\begin{corollary}[The three reciprocal sums]\label{recip}
The reciprocal identities recorded in A322844 are
\begin{align*}
 \sum_{m\geq1}\frac1{a_{2m}}&=2+\frac{\pi^2}{6}-\pi\coth(\pi/2),\\
 \sum_{m\geq1}\frac1{a_{2m+1}}&=15-\frac{3\pi^2}{2},\\
 \sum_{n\geq2}\frac1{a_n}&=17-\frac{4\pi^2}{3}-\pi\coth(\pi/2).
\end{align*}
\end{corollary}
\begin{proof}
All sums converge since $a_n$ grows as a positive constant times $n^4$.
For even indices, $1/a_{2m}=1/m^2-1/(m^2+1/4)$.
The partial-fraction expansion of $\coth z$ \cite{dlmf}, at $z=\pi/2$,
gives $\sum_{m\geq1}(m^2+1/4)^{-1}=\pi\coth(\pi/2)-2$.
For odd $n\geq3$, $1/a_n=12((n^2-1)^{-1}-n^{-2})$;
the two sums are $1/4$ and $\pi^2/8-1$, respectively.
Here the first telescopes after writing $n=2m+1$, and the second follows
from $\sum_{m\geq1}m^{-2}=\pi^2/6$. Adding the parity sums proves the last identity.
\end{proof}

\section{Coverage and verification}
The source statements and their mathematical coverage are as follows.
{\small
\begin{longtable}{@{}p{.43\linewidth}p{.52\linewidth}@{}}
\toprule
Source item&Coverage\\\midrule
\endhead
A322844, COMMENTS, first conjecture; original revision 2, 28 Dec 2018&Proved: $|e_2|=A_n$, and more precisely $e_2=-A_n$, Theorem~\ref{main}.\\
A322844, COMMENTS, second conjecture; original revision 5, 6 Jan 2019&Reproved; prior general rank proof: Vinay (2019), followed by Wicklin (2021). Its specialization gives $e_k=0$ for $k>2$, including $k>n$.\\
Arbitrary row-direction patterns&Not posed; explicit quadratic proved in Corollary~\ref{directions}. Higher-exterior vanishing specializes the same prior rank argument.\\
2012 Miranda construction announcements&Not covered: relation of their unavailable mathematical texts to the spectral coefficient remains unresolved.\\
A317614, determinant; rank link A130130&Reproved: determinant labelled proved by 24 Aug 2018; rank asserted 9 Apr 2023. Neither entry prints an argument.\\
A317614, FORMULA, even-trace decomposition, 28 Dec 2018&Reproved with a spectral interpretation, Corollary~\ref{even}; the minimum-eigenvalue formula is not posed as a conjecture.\\
A317614, concentric-distance matrix, COMMENTS (2021)&Reproved by the centered-absolute-value sum following \eqref{rowcol}.\\
A317614, fourth-power generation algorithm and other arithmetic descriptions&Not covered; these descriptions are separate from the matrix identities.\\
A317297, maximum-eigenvalue conjecture (Spezia, 27 Dec 2018)&Refuted-and-corrected for the literal printed indexing: $26\ne34$ at $r=3$. Proved for $M(2r-2)$, Corollary~\ref{even}.\\
A241016, A317617, sum triangles&Reproved in \eqref{rowcol}; even-order transpose sums proved by De Los Reyes et al. (2007); general formulas recorded in \cite{sums}.\\
A323723, A323724, first off-diagonals&Reproved with prior formulas \cite{diagonals}; parity and the asserted exclusion of terminal digit $4$ proved in Corollary~\ref{lastdigits}.\\
A325516, A325655, all off-diagonals&Proved in \eqref{upper}--\eqref{lower}; prior formulas cited in \cite{diagonals}.\\
A074147, linked antidiagonals&Not asserted in the source: the fixed-size interpretation fails at $n=4$; \eqref{anti} gives the actual entries.\\
A382597, product and trailing zeros of the main antidiagonal&Reproved in Corollary~\ref{antiproduct}; prior statements credited to Spezia (2025).\\
A037270, all-entry sum; A002415&$\sum_{ij}M(n)_{ij}=n^2(n^2+1)/2$ by \eqref{rowcol}; $A_n=\text{A037270}(n)/2$ for even $n$ and $A_n=\text{A002415}(n)$ for odd $n$.\\
A322277, permanent&Not posed as a conjecture; coefficient formula proved in Corollary~\ref{permanent}.\\
A322844, ordinary GFs and parity recurrences&Proved in Corollary~\ref{gf}; prior formulas cited in \cite{posing}.\\
A322844, overall displayed EGF&Refuted-and-corrected under the usual $n!$ convention: $5/24\ne5/2$ at $x^2$; correct formulas \eqref{trueegf}--\eqref{shifted}.\\
A322844, displayed subsequence exponential expressions&Exact factorial conventions proved after \eqref{shifted}; usual $m!$ EGFs in Corollary~\ref{gf}.\\
A322844, three reciprocal sums&Proved in Corollary~\ref{recip}; prior formulas cited in \cite{posing}.\\
A317297, generating functions and recurrence&Proved in Corollary~\ref{triangle}; prior formulas cited in \cite{eigen}.\\
A317297, quartic-difference and ternary-word descriptions (2024--25)&Reproved in Corollary~\ref{triangle}; prior identities credited to Oktar (2024) and Navarrete (2025).\\
A317297, A246697 bisection conjecture; all four A246697 FORMULA conjectures&Proved by the row-sum calculation in Corollary~\ref{triangle}, including $b(r)=c(2r-3)$ for $r\geq2$.\\
A246696, COMMENTS, permutation and parity counts&Permutation property proved; even-row odd count refuted-and-corrected from $[(m+1)/2]$ to $m+1$, Corollary~\ref{permutation}.\\
\bottomrule
\end{longtable}}
The all-dimensional proofs above cover the characteristic polynomial,
spectrum, linear sums, permanent formula and analytic identities.
In the pinned Lean sources,
\href{https://github.com/the-omega-institute/trureturing/blob/d9bb8a406783c5222849b798f3c2ab738f653d75/D5/S3/Quantum/Matrix/SpeziaSerpentineSecondExteriorTrace.lean}{\texttt{result} in the second-exterior file}
proves the formal statement $|e_2(M(n))|=A_n$ for $n>1$, using the coefficient convention
\eqref{exterior}; its \texttt{charpoly\_factor} declaration gives the
$n$-to-$2$ characteristic reduction over $\mathbb Q$ for $n\geq2$.
The companion
\href{https://github.com/the-omega-institute/trureturing/blob/d9bb8a406783c5222849b798f3c2ab738f653d75/D5/S3/Quantum/Matrix/SpeziaSerpentineHigherExteriorTrace.lean}{\texttt{result} in the higher-exterior file}
proves the formal statement $e_k(M(n))=0$ for all $n\geq0$ and $k>2$.
Neither result declaration asserts the spectral or generating-function
corollaries; those are written proofs here. The signed identity $e_2=-A_n$
is also established in the written proof; the first \texttt{result}
declaration asserts its absolute-value version.
The accompanying exact-arithmetic checker compares the full characteristic
polynomial for $0\leq n\leq24$, all principal-minor exterior traces for
$0\leq n\leq8$, the first eight power traces for $2\leq n\leq24$,
and the row, column and diagonal formulas for $1\leq n\leq100$.
It additionally compares the permanent formula with direct permutation sums
through $n=8$, even-dimensional characteristic polynomials through $n=24$,
the explicit $4\times4$ indexing example, and exact generating-function
coefficients and parity recurrences. These finite comparisons complement
the all-$n$ written proofs. Further comparisons cover the first 101 rows
of A246696, the parity counts and disjoint paired parity intervals
through row 100, and first off-diagonal parity and terminal digits
through dimension 100.

\begin{thebibliography}{99}
\bibitem{posing} S. Spezia, OEIS \href{https://oeis.org/A322844}{A322844},
COMMENTS (the two exterior-trace conjectures) and FORMULA;
\href{https://oeis.org/history?seq=A322844\&start=60}{original revisions 2 and 5},
28 December 2018 and 6 January 2019.
\bibitem{matrix} S. Spezia, OEIS \href{https://oeis.org/A317614}{A317614},
1 August 2018, COMMENTS (matrix definition, determinant, and rank link
\href{https://oeis.org/A130130}{A130130}).
The determinant statement is present in revision 89 (24 August 2018);
the rank cross-reference was added in revision 151 (9 April 2023).
\bibitem{sums} OEIS \href{https://oeis.org/A241016}{A241016} and
\href{https://oeis.org/A317617}{A317617}, sum-triangle definitions and formulas;
\href{https://oeis.org/A037270}{A037270}, sum of the first $n^2$ integers;
\href{https://oeis.org/A002415}{A002415}, $(n^4-n^2)/12$.
\bibitem{diagonals} S. Spezia, OEIS
\href{https://oeis.org/A323723}{A323723},
\href{https://oeis.org/A323724}{A323724},
\href{https://oeis.org/A325516}{A325516} and
\href{https://oeis.org/A325655}{A325655}, diagonal-sum definitions and formulas.
\bibitem{antidiagonal} A. Murthy, OEIS
\href{https://oeis.org/A074147}{A074147}, 28 August 2002, defining triangle.
\bibitem{eigen} S. Spezia, OEIS \href{https://oeis.org/A317614}{A317614},
FORMULA, 28 December 2018, even-trace decomposition; O. E. Pol,
OEIS \href{https://oeis.org/A317297}{A317297}, 1 September 2018,
including Spezia's COMMENTS of 27 December 2018,
maximum-eigenvalue conjecture and displayed matrix definition.
\bibitem{permanentref} S. Spezia, OEIS
\href{https://oeis.org/A322277}{A322277}, 1 December 2018,
serpentine-matrix permanent sequence.
\bibitem{dlmf} NIST Digital Library of Mathematical Functions,
\href{https://dlmf.nist.gov/4.36.E3}{Eq.~4.36.3}, partial-fraction expansion
of $\coth z$.
\bibitem{trianglelinks} C. Kimberling, OEIS
\href{https://oeis.org/A246697}{A246697}, 2 September 2014,
FORMULA; \href{https://oeis.org/A246696}{A246696}, defining triangle recurrence.
\bibitem{antiproductref} S. Spezia, OEIS
\href{https://oeis.org/A382597}{A382597}, 31 March 2025,
main-antidiagonal product and trailing-zero statement.
\bibitem{rectangles} J. P. De Los Reyes, A. Das, C. K. Midha and
P. Vellaisamy, \emph{On a method to construct magic rectangles of even order},
Technical Report isid/ms/2007/09, Indian Statistical Institute, Delhi,
5 July 2007, Step A1 and Appendix (i)--(ii).
\href{http://www.isid.ac.in/\string~statmath/eprints/2007/isid200709.pdf}{Official report}.
\bibitem{vinay} M. Vinay, answer to
\href{https://math.stackexchange.com/a/3182097}{``Linear system with constants in arithmetic progression''},
Mathematics Stack Exchange, 10 April 2019 (edited 11 April 2019),
numbered step 2, the general column-span identity.
\bibitem{wicklin} R. Wicklin,
\href{https://blogs.sas.com/content/iml/2021/02/08/matrix-singular-rows-arithmetic-sequences.html}{``A matrix is singular if its rows are arithmetic sequences''},
The DO Loop, SAS, 8 February 2021, ``Why are these matrices singular?''
\bibitem{miranda2012} L. B. Bacelar Miranda,
\href{https://pyaugohy.blogspot.com/2012/09/semi-magic-squares-from-snake-shaped.html}{``Semi-Magic Squares From Snake-Shaped Matrices''},
announcement, 15 September 2012; L'hau\~a Barbosa Pereira de Miranda,
Liuhan Oliveira de Miranda, Lohans de Oliveira Miranda and
Lossian Barbosa Bacelar Miranda,
\href{https://pyaugohy.blogspot.com/2012/09/computation-of-semi-magic-squares.html}{``Computation of Semi-Magic Squares Generated by Serpentine Matrices''},
announcement, 19 September 2012.
\end{thebibliography}
\end{document}
